%
% pi2.tex -- template for standalon tikz images
%
% (c) 2021 Prof Dr Andreas Müller, OST Ostschweizer Fachhochschule
%
\documentclass[tikz]{standalone}
\usepackage{amsmath}
\usepackage{times}
\usepackage{txfonts}
\usepackage{pgfplots}
\usepackage{csvsimple}
\usetikzlibrary{arrows,intersections,math}
\begin{document}
\def\ys{0.5}
\def\skala{0.933}
\pgfmathparse{3.1415*3.1415}
\xdef\piq{\pgfmathresult}
\pgfmathparse{2*3.1415-1.4}
\xdef\lc{\pgfmathresult}
\def\vstep{0.45}

\def\beschriftung#1#2#3#4{
	\draw[color=#3,line width=0.2pt]
		({3.1415},{\ys*(#2(3.1415))})
		--
		(\lc,{\ys*\piq+0.4-(#1)*\vstep});
	\node[color=#3] at (\lc,{\ys*\piq+0.4-(#1)*\vstep}) [right] {$#4$};
	\draw[color=#3,line width=0.7pt]
		plot[domain=-6.3:6.3,samples=300] ({\x},{\ys*#2(\x)});
}

\begin{tikzpicture}[>=latex,thick,scale=\skala,
declare function = {
f0(\X) = 3.2899;
f1(\X) = f0(\X) + 4*(-1)*cos(deg(\X));
f2(\X) = f1(\X) + 4*(0.25)*cos(deg(2*\X));
f3(\X) = f2(\X) + 4*(-0.1111)*cos(deg(3*\X));
f4(\X) = f3(\X) + 4*(0.0625)*cos(deg(4*\X));
f5(\X) = f4(\X) + 4*(-0.04)*cos(deg(5*\X));
f6(\X) = f5(\X) + 4*(0.0277)*cos(deg(6*\X));
f7(\X) = f6(\X) + 4*(-0.02041)*cos(deg(7*\X));
}]

\draw[color=gray,dashed] (3.1415,0) -- ++(0,{\ys*\piq});
\draw[dashed,color=gray] (0,{\ys*\piq}) -- ++(3.1415,0);
\foreach \y in {2,4,6,8}{
	\draw (-0.05,{\ys*\y}) -- ++(0.1,0);
	\node at (0,{\ys*\y}) [left] {$\y$};
}
\draw (-0.05,{\ys*\piq}) -- ++(0.1,0);
\node at (0,{\ys*\piq}) [left] {$\pi^2$};

\begin{scope}
	\clip (-6.28318,-0.2) rectangle (6.28318,{\ys*\piq});
	\draw[color=gray!50,line width=1.2pt]
		plot[domain=-3.1415:3.1415,samples=300] ({\x},{\ys*\x*\x)});

	\begin{scope}[xshift=6.28318cm]
		\draw[color=gray!50,line width=1.2pt]
			plot[domain=-3.1415:3.1415,samples=300]
				({\x},{\ys*\x*\x)});
	\end{scope}

	\begin{scope}[xshift=-6.28318cm]
		\draw[color=gray!50,line width=1.2pt]
			plot[domain=-3.1415:3.1415,samples=300]
				({\x},{\ys*\x*\x)});
	\end{scope}
\end{scope}
\node[color=gray!50] at (-3.1415,{\ys*\piq+0.1}) [right] {$y=x^2$};
\draw (3.1415,{\ys*\piq}) circle[radius=0.1];

\draw[->] (-6.4,0) -- (6.6,0) coordinate[label={$x$}];
\draw[->] (0,-0.4) -- (0,{\ys*11.0}) coordinate[label={right:$y$}];

\beschriftung{7}{f0}{blue!18}{y=f_0(x)\mathstrut}
\beschriftung{6}{f1}{blue!21}{y=f_1(x)\mathstrut}
\beschriftung{5}{f2}{blue!26}{y=f_2(x)\mathstrut}
\beschriftung{4}{f3}{blue!33}{y=f_3(x)\mathstrut}
\beschriftung{3}{f4}{blue!41}{y=f_4(x)\mathstrut}
\beschriftung{2}{f5}{blue!64}{y=f_5(x)\mathstrut}
\beschriftung{1}{f6}{blue!80}{y=f_6(x)\mathstrut}
\beschriftung{0}{f7}{blue!100}{y=f_7(x)\mathstrut}

\def\labelheight{-0.30}
\node at (3.1415,\labelheight) [below] {$\pi\mathstrut$};
\node at ({2*3.1415},\labelheight) [below] {$2\pi\mathstrut$};
\draw (3.1415,-0.05) -- ++(0,0.1);
\draw ({2*3.1415},-0.05) -- ++(0,0.1);
\node at (-3.1415,\labelheight) [below] {$-\pi\mathstrut$};
\node at ({-2*3.1415},\labelheight) [below] {$-2\pi\mathstrut$};
\draw (-3.1415,-0.05) -- ++(0,0.1);
\draw ({-2*3.1415},-0.05) -- ++(0,0.1);
\node at (0,\labelheight) [below] {$0\mathstrut$};

\end{tikzpicture}
\end{document}

